Um die Beziehung $\epsilon^{\prime \alpha \beta} = -\epsilon^{\alpha \beta}$ korrekt herzuleiten, betrachten wir die Transformation eines antisymmetrischen Tensors wie des Levi-Civita-Symbols unter einer linearen Transformation $A$. Diese Transformation ist gegeben durch:

\begin{align*}
\epsilon^{\prime \alpha \beta} &= \det(A) A^\alpha_{\ \gamma} A^\beta_{\ \delta} \epsilon^{\gamma \delta}.
\end{align*}

Da $\epsilon^{\alpha \beta}$ antisymmetrisch ist, gilt $\epsilon^{\alpha \beta} = -\epsilon^{\beta \alpha}$. Für die inverse Transformation $A^{-1}$, die bei orthogonalen Transformationen $A^{-1} = A^T$ ist, ergibt sich:

\begin{align*}
\epsilon^{\alpha \beta} &= \det(A^{-1}) (A^{-1})^\alpha_{\ \gamma} (A^{-1})^\beta_{\ \delta} \epsilon^{\prime \gamma \delta}.
\end{align*}

Da $\det(A^{-1}) = \frac{1}{\det(A)}$, folgt:

\begin{align*}
\epsilon^{\prime \alpha \beta} &= \det(A) A^\alpha_{\ \gamma} A^\beta_{\ \delta} \epsilon^{\gamma \delta}.
\end{align*}

Für orthogonale Transformationen gilt $\det(A) = \pm 1$. Insbesondere für eine Spiegelung ist $\det(A) = -1$, was zu einer Vorzeichenänderung führt:

\begin{align*}
\epsilon^{\prime \alpha \beta} &= - A^\alpha_{\ \gamma} A^\beta_{\ \delta} \epsilon^{\gamma \delta} \quad \text{für eine Spiegelung, da } \det(A) = -1.
\end{align*}

Dies zeigt, dass die transformierte Version des Levi-Civita-Symbols das negative des ursprünglichen Symbols ist, wenn die Transformation eine Spiegelung beinhaltet.

%CHECK: Explizite Darstellung der Vorzeichenänderung bei Spiegelung durch \det(A) = -1.